\section{System Overview}
\label{sec:system}
SHI separates offline corpus preparation from the online retrieval workflow.
During offline preparation, keyframes are enriched with contextual descriptions
generated by a vision-language model (Qwen3-VL-8B-Instruct~\cite{qwen2025qwen3vl}) to
form FrameContext metadata. Together with visual embeddings and spoken
transcripts, these are compiled into dense FAISS indexes~\cite{johnson2019billion}
alongside a lexical BM25 index.
At runtime, as illustrated in Fig.~\ref{fig:layer2}, these artifacts are loaded
as read-only retrieval structures and queried for each event in the active
Query Hypothesis. Corpus-side feature extraction and index construction
therefore remain outside the interactive request path, while online processing
is limited to query encoding, retrieval, score fusion, and temporal alignment.

\begin{figure}[!t]
  \centering
  \includegraphics[width=0.82\textwidth]{layer2.pdf}
  \caption{Multi-event KIS retrieval pipeline in SHI. The Query Hypothesis
  provides an ordered sequence of events. Text events query visual,
  FrameContext, ASR, and lexical indexes, while optional reference images
  contribute visual similarity scores. Component scores are calibrated and
  fused into an event--frame score matrix $S[v,e,f]$ for each candidate video
  $v$. A dynamic-programming decoder then extracts chronologically ordered
  event paths, which are used to rank candidate videos and construct the
  evidence snapshot.}
  \label{fig:layer2}
\end{figure}

\subsection{Multi-event KIS Retrieval Pipeline}
\label{sec:kis-retrieval}

Multi-event KIS is formulated as multimodal event scoring followed by joint
temporal alignment. Given an ordered sequence of events from the Query
Hypothesis, each textual event is evaluated against several retrieval sources:
SigLIP2~\cite{tschannen2025siglip2} provides visual text--frame similarity,
BGE-M3~\cite{chen2024bgem3} provides dense FrameContext and ASR evidence, and
BM25~\cite{robertson2009bm25} provides lexical matching over frame-associated
text. Optional reference images contribute additional
SigLIP2 image--frame similarity scores. The resulting components are calibrated
and combined by event-adaptive weighted fusion into an event--frame score
matrix $S[v,e,f]$ for each candidate video $v$. A dynamic-programming decoder
then selects one canonical frame occurrence for each event subject to strict
temporal ordering. Its objective combines the fused event--frame scores with a
temporal gap penalty that discourages large separations between consecutive
event occurrences. Candidate videos are ranked by the scores of their decoded
paths, so the ranking accounts jointly for event-level evidence and temporal
ordering rather than aggregating independently selected frame matches.

\subsection{Evidence Snapshots and Competition Integration}
\label{sec:evidence-snapshots}

To support post-retrieval interaction without repeating corpus-level search,
SHI captures an in-memory \emph{Evidence Snapshot} after the initial KIS
execution. For the returned candidates, the snapshot retains immutable
event--frame evidence together with the event identifiers, initial decoded
paths, and the decoder configuration used for that search. When an operator
applies an EventTrail action, such as anchoring an accepted occurrence or
rejecting a temporal region, the corresponding constraints are maintained
separately from the frozen evidence and supplied to the temporal decoder. The
decoder then re-computes the feasible path for the selected video using the
retained evidence, without repeating query embedding or corpus index search.
Once the operator selects a frame for submission, SHI's integrated competition
client converts the selection to the required submission format and sends it
to the Distributed Retrieval Evaluation Server
(DRES)~\cite{rossetto2021dres}, reducing the need for manual transfer of frame
coordinates during live search.